\documentclass[11pt]{article}
\usepackage[margin=1in]{geometry}
\usepackage[T1]{fontenc}
\usepackage{lmodern}
\usepackage{url}
\usepackage{setspace}
\setlength{\parindent}{0pt}
\setlength{\parskip}{0.8em}
\onehalfspacing
\begin{document}

\today

Dear Editors,

On behalf of all authors, I am pleased to submit our manuscript entitled ``Self-supervised biomedical signal representation learning from intrapartum cardiotocography for neonatal risk enrichment'' for consideration as an Original Research article in \textit{Frontiers in Signal Processing}, specialty section \textit{Biomedical Signal Processing}.

This manuscript presents a biomedical signal representation-learning study of intrapartum cardiotocography. We model paired fetal heart rate and uterine contraction time series using a compact PatchTST-style self-supervised transformer encoder that combines patch tokenization, masked reconstruction and NT-Xent contrastive consistency. The learned record-level representations are evaluated with classical CTG signal summaries, dynamic phenotype clustering, bootstrap validation, recalibration and external representation/domain-shift assessment.

We believe the manuscript fits \textit{Frontiers in Signal Processing} and the \textit{Biomedical Signal Processing} specialty because its central contribution is a reproducible biomedical signal-processing workflow for long physiological time series, rather than a black-box diagnostic claim. The work asks whether open CTG signals can support reusable representation learning, interpretable trajectory phenotypes and calibrated neonatal risk enrichment under explicit validation boundaries.

The study uses publicly available, de-identified datasets. CTU-UHB/CTU-CHB is used as the primary raw-signal cohort, and CTGDL-FHRMA is used for external representation and domain-shift analysis. Because external neonatal outcome labels were not directly harmonized across datasets, the manuscript deliberately avoids claiming external clinical outcome validation. Instead, it presents CTGDL-FHRMA as a representation-transport and domain-shift test, which we believe is the more valid use of the public data.

The manuscript is original, has not been published previously and is not under consideration elsewhere. All authors have approved the submitted version. The authors declare no competing interests and no external funding. Derived source-data tables, supplementary tables, analysis code and manuscript source files are archived at \url{https://doi.org/10.5281/zenodo.20376424} and \url{https://github.com/Luciky-Leo/ctg-foundation-trajectory}.

Thank you for considering our manuscript.

Sincerely,

Feifan Lu\\
Corresponding Author\\
Department of Obstetrics and Gynecology\\
Changhai Hospital, Naval Medical University\\
Shanghai, China\\
Email: \url{luff94@163.com}

\end{document}
